\subsubsection{BlackBerry OS}

\paragraph{Scheduling Policy and Goals}

BlackBerry~10, through QNX Neutrino, uses priority-based, preemptive
scheduling. Every thread has a priority, and the rule is absolute: if a
higher-priority thread becomes \textsc{ready}, it immediately preempts
all lower-priority threads
\citep{qnx_kernel_sysarch, qnx_sched_policies}. The scheduler's
primary goal is real-time responsiveness and determinism---guaranteeing
that the most important work runs predictably---rather than the
``fairness for everyone'' goal of a general-purpose desktop scheduler.
Priority is the explicit measure (256 levels), and among threads of
equal priority a secondary policy decides order \citep{qnx_kernel_sysarch}.

\paragraph{Scheduling Schemes}

The scheduler picks the highest-priority \textsc{ready} thread. Among
threads sharing the same priority, three policies are available
\emph{per thread} \citep{qnx_kernel_sysarch, qnx_sched_policies}:

\begin{itemize}
\item \textbf{FIFO (SCHED\_FIFO):} the chosen thread runs until it
  blocks or is preempted by a higher-priority thread---it is \emph{not}
  time-sliced against equal-priority peers.
\item \textbf{Round-robin (SCHED\_RR):} identical to FIFO except the
  thread also yields after consuming a timeslice ($4\times$ the clock
  tick; with a CPU faster than 40\,MHz the tick defaults to 1\,ms),
  then goes to the back of its priority queue \citep{qnx_sched_policies}.
\item \textbf{Sporadic (SCHED\_SPORADIC):} gives a thread an execution
  budget within a period, letting its priority oscillate between a
  normal and a low level---designed for Rate Monotonic Analysis of
  mixed periodic/aperiodic real-time work \citep{qnx_kernel_sysarch,
  qnx_sched_policies}.
\end{itemize}

The core data structure is the ready queue, implemented as 256 separate
queues, one per priority level; threads within a queue are normally
ordered FIFO \citep{qnx_kernel_sysarch}. Notably, a server thread
emerging from a \textsc{receive}-blocked state with a client message is
inserted at the \emph{head} of its queue---LIFO---to keep latency low
\citep{qnx_kernel_sysarch}. CPU time is thus allocated either by
blocking (FIFO) or by fixed timeslice (round-robin), and scheduling is
preemptive \citep{qnx_kernel_sysarch, qnx_sched_policies}.

\paragraph{Context Switching}

Context switching is frequent in BlackBerry~10 precisely because so much
OS work happens through message passing between user-space processes---every
Send/Receive/Reply and every interrupt delivery to a driver crosses an
address-space boundary \citep{qnx_arm_memory}. The cost of switching
between two fully private address spaces is higher than switching within
one address space. Older ARM cores reduced this with the FCSE
\citep{qnx_arm_memory}; on the ARMv7-class hardware in BlackBerry~10,
that mechanism is superseded by ASID-tagged TLBs, which serve the same
purpose of avoiding a flush on every switch. The trade-off is the
defining one for the OS: more switching buys responsiveness and
isolation; the overhead is managed by keeping the kernel path short and
the switch cheap.

\paragraph{Process and Thread Handling}

BlackBerry~10 supports the full POSIX model: processes are MMU-protected
from one another, and each process contains one or more threads that
share the process's address space \citep{qnx_kernel_sysarch}. Threads are
kernel-level---the kernel schedules threads directly and individually,
not whole processes, so a multithreaded program genuinely exploits the
scheduler. A thread inherits its scheduling policy from its parent
process but may change it at runtime \citep{qnx_kernel_sysarch,
qnx_sched_policies}. The advantage is true parallelism with fine-grained
scheduling control; the disadvantage is that kernel-level threading
makes every thread operation a (cheap, but real) kernel interaction.

\paragraph{Behaviour Under Different Workloads}

For CPU-bound work, equal-priority threads benefit from round-robin
time-slicing so no one thread starves the others; a single
high-priority compute thread can monopolise the CPU by design, which is
correct for real-time control. For I/O-bound and interactive work, the
strict-priority preemptive model shines: a UI or input thread set at
high priority immediately preempts background computation, giving
snappy, predictable response---exactly what a touch phone needs
\citep{crackberry_qnxhistory}. Long-running batch work is best run at
lower priority or in its own partition (see below) so it consumes only
spare cycles. The OS favours responsiveness for user-facing tasks over
batch throughput.

\paragraph{Multicore and Parallel Processing}

BlackBerry~10 scales across multiple cores and supports SMP (symmetric
multiprocessing) and AMP (asymmetric multiprocessing)
\citep{qnx_smp_multicore}; later QNX releases additionally describe a
QNX-specific refinement called Bound Multiprocessing (BMP), an
improvement on standard SMP processor affinity that lets the designer
bind a thread to a specific core to maximise cache locality while still
allowing the rest of the system to load-balance across cores
\citep{qnx_neutrino_overview}. The trade-off is the classic one: pinning a thread to one core
preserves its warm cache (more cache hits) but can leave that core
overloaded while others idle; pure load balancing keeps cores busy but
sacrifices cache locality. BMP is QNX's way of giving the designer
explicit control over that balance per thread \citep{qnx_neutrino_overview}.

\paragraph{Starvation, Fairness, and Priority Handling}

Pure strict-priority scheduling \emph{can} cause starvation: a stream of
high-priority threads can indefinitely block lower-priority ones. QNX
mitigates this in two principal ways. First, priority inheritance: when
a high-priority client sends a message to a lower-priority server, the
server temporarily inherits the client's priority so it isn't held
back, which also prevents priority inversion
\citep{qnx_ipc, qnx_kernel_sysarch}. Second, QNX's optional adaptive partitioning
scheduler can guarantee each partition a minimum percentage of CPU when
the system is overloaded, so an ``unimportant or untrusted'' partition
cannot monopolise the processor and starve others \citep{qnx_aps}.
Priorities are assigned by the designer and can be adjusted dynamically
(and automatically, under sporadic scheduling)
\citep{qnx_sched_policies}.

\paragraph{Overview --- Processor Management}

Strengths include deterministic real-time response, fine per-thread
control, kernel-level threading, strong multicore support (with BMP
affinity in later releases), and starvation control via priority
inheritance and optional adaptive partitioning
\citep{qnx_kernel_sysarch, qnx_smp_multicore, qnx_neutrino_overview,
qnx_aps}. Weaknesses
are that strict priority demands careful design to avoid starvation and
that heavy message passing makes context switching frequent
\citep{qnx_arm_memory}. In a system mixing interactive apps with
background tasks, QNX performs better because it can guarantee the
interactive partition its CPU share while letting background work
consume only the surplus \citep{qnx_aps}.
